% GENERATED FILE. Edit canonical lessons, then run npm run generate:books.
% canonical-source-hash: fnv1a64:e4985ee405941183
% canonical-lessons: FA-C86-ruye, FA-C86-beyn, FA-C86-samt, FA-C86-taraf, FA-C86-gushe

\chapter{On, Between, Direction, Side, Corner}
\label{ch:fa-on-between-direction-side-corner}
\hlchaptermodality{86}

\begin{chapteropening}
\textbf{By the end of this chapter:} \emph{I can say on, between, a direction, a side and a corner.}

Complete the last lesson of chapter 86: I can say on, between, a direction, a side and a corner.
\end{chapteropening}

\section[\texorpdfstring{\fa{روی} (ruye)}{روی (ruye)}]{\fa{روی} (ruye) — on, on top of}

\label{lesson:FA-C86-ruye}



\begin{quote}
Before the new one: say the Persian for the middle, then the Persian for under.
\end{quote}

\subsection*{You'll want to know: \fa{روی}}
\textbf{\fa{روی}} — \emph{ruye} — \textquotedblleft{}on, on top of\textquotedblright{}.

\begin{cousinweb}
\emph{ruye} is \emph{ru}, \textquotedblleft{}face, surface\textquotedblright{}, with the linking vowel \emph{-ye}: \textquotedblleft{}on the face of\textquotedblright{}. \emph{ru} is native, usually traced to Middle Persian \emph{rōy}, \textquotedblleft{}face\textquotedblright{}. (EtymDB also lists \emph{rōy}, \textquotedblleft{}brass\textquotedblright{}, a different word that looks alike.)
\end{cousinweb}

\begin{grammarlens}[title={What you've built}]
One more word for this chapter.
\end{grammarlens}

\subsection*{Your turn}
\begin{itemize}
\raggedright
  \item \emph{Say it:} \emph{ruye}
  \item \emph{Say it:} \emph{ruye}, once more
  \item \emph{Say it:} say \emph{ruye}
  \item \emph{Recall:} say the Persian for the middle, then the Persian for under, then say \emph{ruye} again
\end{itemize}

\begin{tcolorbox}[breakable,colback=teal!4,colframe=teal!35!black,title={Before you move on}]
What does \fa{روی} mean? (\textquotedblleft{}On, on top of\textquotedblright{}.) Say it once more.
\end{tcolorbox}
\section[\texorpdfstring{\fa{بین} (beyn)}{بین (beyn)}]{\fa{بین} (beyn) — between}

\label{lesson:FA-C86-beyn}



\begin{quote}
Before the new one: say the Persian for under, then the Persian for on.
\end{quote}

\subsection*{You'll want to know: \fa{بین}}
\textbf{\fa{بین}} — \emph{beyn} — \textquotedblleft{}between\textquotedblright{}.

\begin{cousinweb}
\emph{beyn} is borrowed from Arabic \emph{bayna}, \textquotedblleft{}between\textquotedblright{}. Its Arabic root is \emph{b-y-n} (Lane).
\end{cousinweb}

\begin{grammarlens}[title={What you've built}]
One more word for this chapter.
\end{grammarlens}

\subsection*{Your turn}
\begin{itemize}
\raggedright
  \item \emph{Say it:} \emph{beyn}
  \item \emph{Say it:} \emph{beyn}, once more
  \item \emph{Say it:} say \emph{beyn}
  \item \emph{Recall:} say the Persian for under, then the Persian for on, then say \emph{beyn} again
\end{itemize}

\begin{tcolorbox}[breakable,colback=teal!4,colframe=teal!35!black,title={Before you move on}]
What does \fa{بین} mean? (\textquotedblleft{}Between\textquotedblright{}.) Say it once more.
\end{tcolorbox}
\section[\texorpdfstring{\fa{سمت} (samt)}{سمت (samt)}]{\fa{سمت} (samt) — a direction}

\label{lesson:FA-C86-samt}



\begin{quote}
Before the new one: say the Persian for on, then the Persian for between.
\end{quote}

\subsection*{You'll want to know: \fa{سمت}}
\textbf{\fa{سمت}} — \emph{samt} — \textquotedblleft{}a direction\textquotedblright{}.

\begin{cousinweb}
\emph{samt} is borrowed from Arabic \emph{samt}, \textquotedblleft{}way, direction\textquotedblright{}. Its Arabic root is \emph{s-m-t} (Lane). English \emph{zenith} comes from the Arabic phrase \emph{samt ar-raʾs}, \textquotedblleft{}direction of the head\textquotedblright{} (EtymDB).
\end{cousinweb}

\begin{grammarlens}[title={What you've built}]
One more word for this chapter.
\end{grammarlens}

\subsection*{Your turn}
\begin{itemize}
\raggedright
  \item \emph{Say it:} \emph{samt}
  \item \emph{Say it:} \emph{samt}, once more
  \item \emph{Say it:} say \emph{samt}
  \item \emph{Recall:} say the Persian for on, then the Persian for between, then say \emph{samt} again
\end{itemize}

\begin{tcolorbox}[breakable,colback=teal!4,colframe=teal!35!black,title={Before you move on}]
What does \fa{سمت} mean? (\textquotedblleft{}A direction\textquotedblright{}.) Say it once more.
\end{tcolorbox}
\section[\texorpdfstring{\fa{طرف} (taraf)}{طرف (taraf)}]{\fa{طرف} (taraf) — a side}

\label{lesson:FA-C86-taraf}



\begin{quote}
Before the new one: say the Persian for between, then the Persian for a direction.
\end{quote}

\subsection*{You'll want to know: \fa{طرف}}
\textbf{\fa{طرف}} — \emph{taraf} — \textquotedblleft{}a side\textquotedblright{}.

\begin{cousinweb}
\emph{taraf} is borrowed from Arabic \emph{ṭaraf}, \textquotedblleft{}edge, side\textquotedblright{}. Its Arabic root is \emph{ṭ-r-f} (Lane).
\end{cousinweb}

\begin{grammarlens}[title={What you've built}]
One more word for this chapter.
\end{grammarlens}

\subsection*{Your turn}
\begin{itemize}
\raggedright
  \item \emph{Say it:} \emph{taraf}
  \item \emph{Say it:} \emph{taraf}, once more
  \item \emph{Say it:} say \emph{taraf}
  \item \emph{Recall:} say the Persian for between, then the Persian for a direction, then say \emph{taraf} again
\end{itemize}

\begin{tcolorbox}[breakable,colback=teal!4,colframe=teal!35!black,title={Before you move on}]
What does \fa{طرف} mean? (\textquotedblleft{}A side\textquotedblright{}.) Say it once more.
\end{tcolorbox}
\section[\texorpdfstring{\fa{گوشه} (gushe)}{گوشه (gushe)}]{\fa{گوشه} (gushe) — a corner}

\label{lesson:FA-C86-gushe}



\begin{quote}
Before the new one: say the Persian for a direction, then the Persian for a side.
\end{quote}

\subsection*{You'll want to know: \fa{گوشه}}
\textbf{\fa{گوشه}} — \emph{gushe} — \textquotedblleft{}a corner\textquotedblright{}.

\begin{cousinweb}
\emph{gushe} is a native Persian word, inherited rather than borrowed: it continues Middle Persian \emph{gōšag}, \textquotedblleft{}corner\textquotedblright{} (EtymDB).
\end{cousinweb}

\begin{grammarlens}[title={What you've built}]
One more word for this chapter.
\end{grammarlens}

\subsection*{Your turn}
\begin{itemize}
\raggedright
  \item \emph{Say it:} \emph{gushe}
  \item \emph{Say it:} \emph{gushe}, once more
  \item \emph{Say it:} say \emph{gushe}
  \item \emph{Recall:} say the Persian for a direction, then the Persian for a side, then say \emph{gushe} again
\end{itemize}

\begin{tcolorbox}[breakable,colback=teal!4,colframe=teal!35!black,title={Before you move on}]
What does \fa{گوشه} mean? (\textquotedblleft{}A corner\textquotedblright{}.) Say it once more.
\end{tcolorbox}
